\documentclass{article}
% \usepackage[legalpaper,margin=0.75in]{geometry}
\usepackage[a4paper, margin=1in]{geometry}

\usepackage[utf8]{inputenc}
\usepackage[T1]{fontenc}
% \usepackage[default]{comicneue}

\usepackage{hyperref}
\hypersetup{
	colorlinks=true, %set true if you want colored links
	linktoc=all,     %set to all if you want both sections and subsections linked
	linkcolor=black,  %choose some color if you want links to stand out
}

% \usepackage{url}

\usepackage[style=alphabetic]{biblatex}
\addbibresource{bibliography.bib}

\usepackage{graphicx}
\usepackage{enumitem}
\usepackage{booktabs}
\usepackage{csquotes}
% \usepackage{emptypage}
\usepackage{subcaption}
\usepackage{multicol}
\usepackage[dvipsnames]{xcolor}
% \usepackage[usenames, dvipsnames]{xcolor}
\usepackage{amsmath, mathtools, amssymb} 

% \usepackag{amsfonts}
% text and math fonts
\usepackage{newpxtext} 
\usepackage[euler-digits]{eulervm}

\usepackage{amsthm}
\usepackage{mathrsfs}
\usepackage{cancel}
\usepackage{bm}
\usepackage{systeme}
\usepackage{stmaryrd}
\usepackage{dsfont}



\DeclareCiteCommand{\cite}
{\usebibmacro{prenote}}
{\usebibmacro{citeindex}%
	\printtext[bibhyperref]{\usebibmacro{cite}}}
{\multicitedelim}
{\usebibmacro{postnote}}

\DeclareCiteCommand*{\cite}
{\usebibmacro{prenote}}
{\usebibmacro{citeindex}%
	\printtext[bibhyperref]{\usebibmacro{citeyear}}}
{\multicitedelim}
{\usebibmacro{postnote}}

\DeclareCiteCommand{\parencite}[\mkbibparens]
{\usebibmacro{prenote}}
{\usebibmacro{citeindex}%
	\printtext[bibhyperref]{\usebibmacro{cite}}}
{\multicitedelim}
{\usebibmacro{postnote}}

\DeclareCiteCommand*{\parencite}[\mkbibparens]
{\usebibmacro{prenote}}
{\usebibmacro{citeindex}%
	\printtext[bibhyperref]{\usebibmacro{citeyear}}}
{\multicitedelim}
{\usebibmacro{postnote}}

\DeclareCiteCommand{\footcite}[\mkbibfootnote]
{\usebibmacro{prenote}}
{\usebibmacro{citeindex}%
	\printtext[bibhyperref]{ \usebibmacro{cite}}}
{\multicitedelim}
{\usebibmacro{postnote}}

\DeclareCiteCommand{\footcitetext}[\mkbibfootnotetext]
{\usebibmacro{prenote}}
{\usebibmacro{citeindex}%
	\printtext[bibhyperref]{\usebibmacro{cite}}}
{\multicitedelim}
{\usebibmacro{postnote}}

\DeclareCiteCommand{\textcite}
{\boolfalse{cbx:parens}}
{\usebibmacro{citeindex}%
	\printtext[bibhyperref]{\usebibmacro{textcite}}}
{\ifbool{cbx:parens}
	{\bibcloseparen\global\boolfalse{cbx:parens}}
	{}%
	\multicitedelim}
{\usebibmacro{textcite:postnote}}



\let\implies\Rightarrow
\let\impliedby\Leftarrow
\let\iff\Leftrightarrow
\let\epsilon\varepsilon
\newcommand\contra{\scalebox{1.1}{$\lightning$}}
\let\emptyset\varnothing


\definecolor{correct}{HTML}{009900}
\newcommand\correct[2]{\ensuremath{\:}{\color{red}{#1}}\ensuremath{\to }{\color{correct}{#2}}\ensuremath{\:}}
\newcommand\green[1]{{\color{correct}{#1}}}


\newcommand\hr{
	\noindent\rule[0.5ex]{\linewidth}{0.5pt}
}


% hide parts
\newcommand\hide[1]{}



% si unitx
\usepackage{siunitx}
\sisetup{locale = FR}
% \renewcommand\vec[1]{\mathbf{#1}}
\newcommand\mat[1]{\mathbf{#1}}


% tikz
\usepackage{tikz}
\usepackage{tikz-cd}
\usetikzlibrary{intersections, angles, quotes, calc, positioning}
\usetikzlibrary{arrows.meta}
\usepackage{pgfplots}
\pgfplotsset{compat=1.13}



\tikzset{
	force/.style={thick, {Circle[length=2pt]}-stealth, shorten <=-1pt}
}

% theorems
\makeatother
\usepackage{thmtools}
\usepackage[framemethod=TikZ]{mdframed}
\mdfsetup{skipabove=1em,skipbelow=0em}


% \theoremstyle{plain}

\declaretheoremstyle[
	headfont=\bfseries\sffamily\color{ForestGreen!70!black}, bodyfont=\normalfont,
	mdframed={
			linewidth=2pt,
			rightline=false, topline=false, bottomline=false,
			linecolor=ForestGreen, backgroundcolor=ForestGreen!5,
		}
]{thmgreenbox}

\declaretheoremstyle[
	headfont=\bfseries\sffamily\color{NavyBlue!70!black}, bodyfont=\normalfont,
	mdframed={
			linewidth=2pt,
			rightline=false, topline=false, bottomline=false,
			linecolor=NavyBlue, backgroundcolor=NavyBlue!5,
		}
]{thmbluebox}

\declaretheoremstyle[
	headfont=\bfseries\sffamily\color{RoyalPurple}, bodyfont=\normalfont\itshape,
	mdframed={
			linewidth=2pt,
			rightline=false, topline=false, bottomline=false,
			linecolor=Rhodamine, backgroundcolor=Rhodamine!5,
		}
]{thmpinkbox}

\declaretheoremstyle[
	headfont=\bfseries\sffamily\color{NavyBlue!70!black}, bodyfont=\normalfont,
	mdframed={
			linewidth=2pt,
			rightline=false, topline=false, bottomline=false,
			linecolor=NavyBlue
		}
]{thmblueline}

\declaretheoremstyle[
	headfont=\bfseries\sffamily\color{RawSienna!70!black}, bodyfont=\normalfont\itshape,
	mdframed={
			linewidth=2pt,
			rightline=false, topline=false, bottomline=false,
			linecolor=RawSienna, backgroundcolor=RawSienna!5,
		}
]{thmredbox}

\declaretheoremstyle[
	headfont=\bfseries\sffamily\color{RawSienna!70!black}, bodyfont=\normalfont,
	numbered=no,
	mdframed={
			linewidth=2pt,
			rightline=false, topline=false, bottomline=false,
			linecolor=RawSienna, backgroundcolor=RawSienna!1,
		},
	qed=\qedsymbol
]{thmproofbox}

\declaretheoremstyle[
	headfont=\bfseries\sffamily\color{NavyBlue!70!black}, bodyfont=\normalfont,
	numbered=no,
	mdframed={
			linewidth=2pt,
			rightline=false, topline=false, bottomline=false,
			linecolor=NavyBlue, backgroundcolor=NavyBlue!1,
		},
]{thmexplanationbox}

\declaretheorem[style=thmgreenbox, name=Definition, numberwithin=section]{definition}
\declaretheorem[style=thmbluebox, name=Example, numberwithin=section]{eg}
\declaretheorem[style=thmpinkbox, name=Exercise, numberwithin=section]{ex}
\declaretheorem[style=thmredbox, name=Proposition, numberwithin=section]{prop}
\declaretheorem[style=thmredbox, name=Theorem, numberwithin=section]{theorem}
\declaretheorem[style=thmredbox, name=Lemma, numberwithin=section]{lemma}
\declaretheorem[style=thmredbox, numbered=no, name=Corollary, numberwithin=section]{corollary}

\declaretheorem[style=thmproofbox, name=Proof]{replacementproof}
\renewenvironment{proof}[1][\proofname]{\vspace{-10pt}\begin{replacementproof}}{\end{replacementproof}}


\declaretheorem[style=thmexplanationbox, name=Proof]{tmpexplanation}
\newenvironment{explanation}[1][]{\vspace{-10pt}\begin{tmpexplanation}}{\end{tmpexplanation}}

\declaretheorem[style=thmblueline, numbered=no, name=Remark]{remark}
\declaretheorem[style=thmblueline, numbered=no, name=Note]{note}

\newtheorem*{uovt}{UOVT}
\newtheorem*{notation}{Notation}
\newtheorem*{previouslyseen}{As previously seen}
\newtheorem*{problem}{Problem}
\newtheorem*{observe}{Observe}
\newtheorem*{property}{Property}
\newtheorem*{intuition}{Intuition}


\usepackage{etoolbox}
\AtEndEnvironment{vb}{\null\hfill$\diamond$}%
\AtEndEnvironment{intermezzo}{\null\hfill$\diamond$}%
% \AtEndEnvironment{opmerking}{\null\hfill$\diamond$}%

% http://tex.stackexchange.com/questions/22119/how-can-i-change-the-spacing-before-theorems-with-amsthm
\makeatletter
% \def\thm@space@setup{%
%   \thm@preskip=\parskip \thm@postskip=0pt
% }

\newcommand{\oefening}[1]{%
	\def\@oefening{#1}%
	\subsection*{Oefening #1}
}

\newcommand{\suboefening}[1]{%
	\subsubsection*{Oefening \@oefening.#1}
}

\newcommand{\exercise}[1]{%
	\def\@exercise{#1}%
	\subsection*{Exercise #1}
}

\newcommand{\subexercise}[1]{%
	\subsubsection*{Exercise \@exercise.#1}
}


\usepackage{xifthen}

\def\testdateparts#1{\dateparts#1\relax}
\def\dateparts#1 #2 #3 #4 #5\relax{
	\marginpar{\small\textsf{\mbox{#1 #2 #3 #5}}}
}

\def\@lesson{}%
\newcommand{\lesson}[3]{
	\ifthenelse{\isempty{#3}}{%
		\def\@lesson{Lecture #1}%
	}{%
		\def\@lesson{Lecture #1: #3}%
	}%
	\subsection*{\@lesson}
	\testdateparts{#2}
}

% \renewcommand\date[1]{\marginpar{#1}}


% % fancy headers
%
% \usepackage{fancyhdr}
% \pagestyle{fancy}
% % \fancyhead[LE,RO]{Gilles Castel}
% % \fancyhead[RO,LE]{\@lesson}
% % \fancyhead[RE,LO]{}
% % \fancyfoot[LE,RO]{\thepage}
% % \fancyfoot[C]{\leftmark}
% % \fancyhf{}
% \fancyhf[roh,leh]{\thepage}
% \fancyhf[reh]{\nouppercase{\leftmark}}
% \fancyhf[loh]{\nouppercase{\rightmark}}
% \cfoot{}

% \setlength{\headheight}{15pt}
% \renewcommand{\headrulewidth}{0pt}
% \renewcommand{\footrulewidth}{0pt}


\makeatother

% notes
% \usepackage{todonotes}
\usepackage{tcolorbox}

\tcbuselibrary{breakable}
\newenvironment{verbetering}{\begin{tcolorbox}[
			arc=0mm,
			colback=white,
			colframe=green!60!black,
			title=Opmerking,
			fonttitle=\sffamily,
			breakable
		]}{\end{tcolorbox}}

\newenvironment{noot}[1]{\begin{tcolorbox}[
			arc=0mm,
			colback=white,
			colframe=white!60!black,
			title=#1,
			fonttitle=\sffamily,
			breakable
		]}{\end{tcolorbox}}


% figure support
\usepackage{import}
\usepackage{xifthen}
\pdfminorversion=7
\usepackage{pdfpages}
\usepackage{transparent}
\newcommand{\incfig}[1]{%
	\def\svgwidth{\columnwidth}
	\import{./figures/}{#1.pdf_tex}
}

% %http://tex.stackexchange.com/questions/76273/multiple-pdfs-with-page-group-included-in-a-single-page-warning
\pdfsuppresswarningpagegroup=1


\newcommand{\one}{{\mathds 1}}
\newcommand{\calA}{{\mathcal A}}
\newcommand{\calB}{{\mathcal B}}
\newcommand{\relation}{{\mathcal R}}
\newcommand{\sigmaalg}{{\mathcal M}}
\newcommand{\calF}{{\mathcal F}}
\newcommand{\calR}{{\mathcal R}}
\newcommand{\calE}{{\mathcal E}}

\newcommand\restr[2]{{% we make the whole thing an ordinary symbol
\left.\kern-\nulldelimiterspace % automatically resize the bar with \right
#1 % the function
\littletaller % pretend it's a little taller at normal size
\right|_{#2} % this is the delimiter
}}

\newcommand{\littletaller}{\mathchoice{\vphantom{\big|}}{}{}{}}

\newcommand{\outermeasure}{{{\mu^{\star}}}}

\def\morphism#1{\overset{#1}\longrightarrow}
\def\xto#1{\xrightarrow{#1}}
\def\sf#1{\mathsf{#1}}
\def\mbb#1{\mathbb{#1}}
\def\mfk#1{\mathfrak{#1}}
\def\bN{\mbb{N}}
\def\bC{\mbb{C}}
\def\bR{\mbb{R}}
\def\bQ{\mbb{Q}}
\def\bZ{\mbb{Z}}
\def\bB{\mbb{B}}
\def\catC{\sf{C}}
\def\HomC{\sf{Hom}_{\catC}}
\def\Hom{\sf{Hom}}
\def\Aut{\sf{Aut}}
\def\End{\sf{End}}
\def\Obj{\sf{Obj}}
\newcommand\quotient[2]{{{\displaystyle #1}}/{{\displaystyle #2}}}

% gonna keep these commands for backwards compatibility
\newcommand{\N}{{\mathbb N}}
\newcommand{\R}{{\mathbb R}}
\newcommand{\powerset}{{\mathscr P}}
\newcommand{\inv}{^{\raisebox{.2ex}{$\scriptscriptstyle-1$}}}

\DeclareMathOperator{\image}{im}
\DeclareMathOperator{\Id}{id}
% \DeclareMathOperator{\Hom}{Hom}
% \DeclareMathOperator{\End}{End}
% \DeclareMathOperator{\Aut}{Aut}
\DeclareMathOperator{\Span}{span}





\author{pika}
\date{\today}


\title{Aluffi Chapter-1 Solutions}
\begin{document}
\maketitle
% \tableofcontents

\begin{ex}\label{ex:rmodziscirclegroup}
  \cite[Ex-1.6]{aluffi_algebra_2009}
  Define a relation \(\sim\) on \(\bR\) as \(a \sim b\) if \(a-b\in \bZ\). 
  Show that \(\sim\) is an equivalence relation on \(\bR\). Find a \emph{compelling} 
  description for \(\bR/\sim\). Do the same for the relation \(\equiv\) 
  on the plane \(\bR\times \bR\) defined by \((a_{1}, a_{2}) \equiv (b_{1}, b_{2}) \iff b_{1}-a_{1} \in \bZ\) and \(b_{2}-a_{2}\in \bZ\).
\end{ex}
\begin{proof}
  Reflexivity holds because \(x-x=0\in \bZ\).
  Symmetry is due to \(y-x = -(x-y)\in \bZ\).
  Finally, for \(x,y,z \in \bR\) with \(x - y = a \in \bZ\) and \(y - z = b \in \bZ\), 
  we can write \(x - z = x - y - (z - y) = a + b \in \bZ\) which establishes 
  transitivity. This completes the porof that \(\sim\) is an equivalence relation. 

  Next, we will show that \(\quotient{\bR}{\sim}\) is in bijection to the following set: 
  \[
    S^{1} \coloneqq \left\{x \in \bC; \|x\| = 1\right\}
  \]

  For this we introduce a map 
  \[
    g : \bR \to S^{1}
  \]
  defined by
  \[
    x \mapsto e^{2\pi i x}
  \]
  That \(g\) is well defined is obvious from intro analysis.
  We know from intro analysis that \(g\) is a surjection. 

  Then, we shall invoke \cite[Thm-2.7]{aluffi_algebra_2009} on the diagram
  \[
    \begin{tikzcd}
      \bR && S^{1} \\
      \\
      {\quotient{\bR}{\sim}}
      \arrow["g", from=1-1, to=1-3]
      \arrow["p"', from=1-1, to=3-1]
      \arrow["f"', from=3-1, to=1-3]
    \end{tikzcd}
  \]
  This gives us a well defined injection \(f:\quotient{\bR}{\sim}\to S^{1}\). 
  We get that \(\image f = \image g\) and since \(g\) is surjective \(\image f = S^{1}\)
  which completes the proof.

  Next, we make an important observation. 
  \begin{observe}
    Since any pair of points in any set in \(\quotient{\bR}{\sim}\) have the same 
    fractional part, quotienting \(\bR\) by \(\sim\) has the effect of 
    treating the integral part of any real as insignificant and of course, those reals with 
    identical fractional parts are identified as equivalent. 
    Motivated by this, hereafter, we will simply use the notation \(\quotient{\bR}{\bZ}\)
    for \(\quotient{\bR}{\sim}\).
  \end{observe}

  Since \(\quotient{\bR}{\bZ} \cong S^{1}\) we are tempted to believe that
  \(\quotient{\bR^{2}}{\bZ^{2}}\) is in bijection with \(S^{1}\times S^{1}\).
  We shall prove this. 

  First, notice that there is a bijection \(f: \quotient{\bR}{\bZ} \to S^{1}\). 
  This gives us a bijection

  \[
    \quotient{\bR}{\bZ} \times \quotient{\bR}{\bZ} \to S^{1} \times S^{1}
  \]

  by the rule

  \[
    (x + \bZ, y + \bZ) \mapsto (f(x), f(y))
  \]

  So, we simply need to prove that there is a bijection between 
  \(\quotient{\bR^{2}}{\bZ^{2}}\) and \(\quotient{\bR}{\bZ}\times\quotient{\bR}{\bZ}\).

  Here, we shall use the universal property of products as in 
  \cite[Prop-1.3.1]{kim_rough_2018}.

  For this, consider the function

  \[
    f : \quotient{\bR^{2}}{\bZ^{2}} \to \quotient{\bR}{\bZ}
  \]

  defined by 

  \[
    (x,y) + \bZ \times \bZ \mapsto x + \bZ
  \]

  and the function 

  \[
    g: \quotient{\bR^{2}}{\bZ^{2}} \to \quotient{\bR}{\bZ} 
  \]

  defined by 

  \[
    (x,y) + \bZ \times \bZ \mapsto y + \bZ
  \]

  Next, consider the diagram 

  % https://q.uiver.app/#q=WzAsNixbMiwyLCJcXG1hdGhiYntSfS9cXG1hdGhiYntafSBcXHRpbWVzIFxcbWF0aGJie1J9L1xcbWF0aGJie1p9Il0sWzIsMCwiXFxtYXRoYmJ7Un1eezJ9L1xcbWF0aGJie1p9XnsyfSJdLFs0LDYsIlNeMSBcXHRpbWVzIFNeMSJdLFszLDYsIlxcYnVsbGV0Il0sWzAsMywiXFxtYXRoYmJ7Un0vXFxtYXRoYmJ7Wn0iXSxbNCwzLCJcXG1hdGhiYntSfS9cXG1hdGhiYntafSJdLFsxLDBdLFsyLDNdLFsxLDQsImYiLDJdLFswLDQsIlxccGlfMSJdLFsxLDUsImciXSxbMCw1LCJcXHBpXzIiLDJdXQ==
  \[
    \begin{tikzcd}
      && {\mathbb{R}^{2}/\mathbb{Z}^{2}} && \\
      \\
      && {\mathbb{R}/\mathbb{Z} \times \mathbb{R}/\mathbb{Z}} \\
      {\mathbb{R}/\mathbb{Z}} &&&& {\mathbb{R}/\mathbb{Z}} \\
      \\
      \arrow[dashed, "h"', from=1-3, to=3-3]
      \arrow["f"', from=1-3, to=4-1]
      \arrow["g", from=1-3, to=4-5]
      \arrow["{\pi_1}", from=3-3, to=4-1]
      \arrow["{\pi_2}"', from=3-3, to=4-5]
    \end{tikzcd}
  \]

  By universal properties of products, there is a map 

  \[
    h: \quotient{\bR^{2}}{\bZ^{2}} \to \quotient{\bR}{\bZ}\times \quotient{\bR}{\bZ}
  \]

  which satisfies 

  \[
    f = \pi_{1}\circ h, \qquad g = \pi_{2}\circ h
  \]

  It is obvious that \(f\) and \(g\) are surjective and from this follows the surjectivity of \(h\).

  % We argue that \(h\) is a surjection. Indeed, take any 
  % \((x + \bZ, y + \bZ) \in \quotient{\bR}{\bZ}\times \quotient{\bR}{\bZ}\)
  %
  % For this, take \((x,y) + \bZ\times \bZ \in \quotient{\bR^{2}}{\bZ^{2}}\) and notice that 
  % \(f((x,y) + \bZ\times \bZ) = x + \bZ\) and \(g((x,y) + \bZ \times \bZ) = y + \bZ\)
  %
  % This tells us that \((\pi_1 \circ h)((x,y)+\bZ\times \bZ) = x + \bZ\) 
  % and \((\pi_2 \circ h)((x,y)+\bZ \times \bZ) = y + \bZ\).
  %
  % From this it is obvious that \(h((x,y) + \bZ\times \bZ) = (x+\bZ, y+\bZ)\) so, \(h\) 
  % is surjective. 

  Next, take any 

  \[
    (x,y) + \bZ^{2} \in \quotient{\bR^{2}}{\bZ^{2}}
  \]

  and 

  \[
    (x', y') + \bZ^{2} \in \quotient{\bR^{2}}{\bZ^{2}}
  \]

  so that \(h\left((x,y) + \bZ^{2}\right) = h\left((x',y') + \bZ^{2}\right)\)

  Then, \((\pi_1\circ h) ((x,y) + \bZ^{2}) = x+\bZ\) and \((\pi_2 \circ h) ((x,y)+\bZ^{2}) = y+\bZ\)
  and a similar result holds for \((x',y')+\bZ^{2}\).

  So, if 

  \[
    (x+\bZ, y+\bZ) = (x'+\bZ, y'+\bZ)
  \]
  
  then we have 

  \[
    x+\bZ = x'+\bZ, \qquad y+\bZ = y'+\bZ 
  \]

  This tells us that \(x \equiv x' \pmod{\bZ}\) and \(y \equiv y' \pmod{\bZ}\). This is 
  precisely the condition for 

  \[
    (x,y) \equiv (x', y') \pmod{\bZ^{2}}
  \]

  which is another way of saying that \((x,y)+\bZ^{2} = (x',y')+\bZ^{2}\) which 
  proves that \(h\) is an injection thereby completing the proof.
\end{proof}

\begin{remark} 
  It is worth mentioning that the second half of \ref{ex:rmodziscirclegroup}
  could be proved by invoking the universal property of the quotient \(\quotient{\bR^{2}}{\bZ^{2}}\)
  instead of invoking the universal property of the product \(\quotient{\bR}{\bZ}\times \quotient{\bR}{\bZ}\).

  In fact, we can combine the two universal properties together. We can use 
  the property on \(\quotient{\bR^{2}}{\bZ^{2}}\) to find a unique map \(s\) that 
  is injective and realize that \(s\) is the same as \(h\) by invoking 
  both universal properties.


  % https://q.uiver.app/#q=WzAsNixbMiwyLCJcXG1hdGhiYntSfS9cXG1hdGhiYntafSBcXHRpbWVzIFxcbWF0aGJie1J9L1xcbWF0aGJie1p9Il0sWzIsMCwiXFxtYXRoYmJ7Un1eezJ9L1xcbWF0aGJie1p9XnsyfSJdLFszLDYsIlxcYnVsbGV0Il0sWzAsMywiXFxtYXRoYmJ7Un0vXFxtYXRoYmJ7Wn0iXSxbNCwzLCJcXG1hdGhiYntSfS9cXG1hdGhiYntafSJdLFswLDAsIlxcbWF0aGJie1J9XnsyfSJdLFsxLDAsImg9cyIsMCx7InN0eWxlIjp7ImJvZHkiOnsibmFtZSI6ImRhc2hlZCJ9fX1dLFsxLDMsImYiLDJdLFswLDMsIlxccGlfMSJdLFsxLDQsImciXSxbMCw0LCJcXHBpXzIiLDJdLFs1LDEsInAiXSxbNSwwLCJyIl1d
  \[\begin{tikzcd}
    {\mathbb{R}^{2}} && {\mathbb{R}^{2}/\mathbb{Z}^{2}} && \\
    \\
    && {\mathbb{R}/\mathbb{Z} \times \mathbb{R}/\mathbb{Z}} \\
    {\mathbb{R}/\mathbb{Z}} &&&& {\mathbb{R}/\mathbb{Z}} \\
    \arrow["p", from=1-1, to=1-3]
    \arrow["r", from=1-1, to=3-3]
    \arrow["{h,s}", dashed, from=1-3, to=3-3]
    \arrow["f"', from=1-3, to=4-1]
    \arrow["g", from=1-3, to=4-5]
    \arrow["{\pi_1}", from=3-3, to=4-1]
    \arrow["{\pi_2}"', from=3-3, to=4-5]
  \end{tikzcd}\]
  
  Here, \(p\) is the projection of \(\mathbb{R}^{2}\) onto the quotient \(\quotient{\bR^{2}}{\bZ^{2}}\),
  \(r:\bR^{2}\to \quotient{\bR}{\bZ} \times \quotient{\bR}{\bZ}\) is the map 
  \((x,y)\mapsto (x+\bZ, y+\bZ)\).
  
\end{remark}

\begin{ex}\cite[Prop-2.1.2]{aluffi_algebra_2009}
	Let \(f:A\to B\) be a function between any two sets \(A\) and \(B\).
	Prove that \(f\) has a right inverse if and only if \(f\) is surjective.
\end{ex}
\begin{proof}
  (\(\Longrightarrow\)) Suppose \(\image f = B\). Then, for any \(b \in B\) 
  the set \(f^{-1}\left(\{b\}\right)\) (the fiber of \(f\) at \(b\)) is nonempty. 
  So, by axiom of choice, for every \(b \in B\) we choose a \(b^{\dagger}\in A\) so that 
  \(b^{\dagger} \in f^{-1}\left(\{b\}\right)\). By this, we define the map

  \begin{align*}
    g \colon B &\longrightarrow A \\
    b &\longmapsto b^{\dagger} \in f^{-1}(\{b\})
  \end{align*}

  Next, take any \(b \in b\). Then, \((f\circ g)(b) = f(g(b)) = f(b^{\dagger}) = b\). Therefore, 
  \(f \circ g = \one_{B}\).

  (\(\Longleftarrow\))
  Suppose there is a map \(g: B \to A\) so that \(f \circ g = \one_{B}\).
  Then, pick any \(b \in B\). Then, \(g(b)\) is an element of \(A\) so that 

  \begin{equation}
    f(g(b)) = b
  \end{equation}

  proving that \(\image f = B\).

\end{proof}
\begin{remark}
  Notice that we require no further condition on the sets \(A\) and \(B\). 
\end{remark}

\begin{ex}\cite[Ex-2.6]{aluffi_algebra_2009}
  Explain how any function \(f: A \rightarrow B\) determines a section of the projection 
  \(\pi_{A}: A\times B \to A\).
\end{ex}
\begin{proof}

  The map \(g : A \to A \times B\) defined by \(a \mapsto (a, f(a))\)
  is a section of \(\pi_{A}\). Indeed, given any \(a\in A\), 

  \begin{align*}
    (\pi_{A}\circ g)(a) &= \pi_{A}(g(a))\\
                        &= \pi_{A}((a,f(a)) \\
                        &= a
  \end{align*}
  
  Therefore, \(\pi\circ g = \one_{A}\).

\end{proof}

\begin{definition}\label{def:epimorphism}
  A function \(f: A\to B\) between any two sets \(A\) and \(B\) is said to be an 
  \emph{epimorphism} if for every set \(Z\) and a pair of maps 
  \(\alpha',~\alpha'': B \to Z\) it holds that 
  \(\alpha' \circ f = \alpha'' \circ f\) then \(\alpha' = \alpha''\).
\end{definition}
\begin{ex}\cite[Ex-2.5]{aluffi_algebra_2009}
  A function \(f: A \to B\) between any two sets \(A\) and \(B\) is surjective if and 
  only if it is an epimorphism (\ref{def:epimorphism}).
\end{ex}
\begin{proof}
  \((\Longrightarrow)\) Suppose \(f: A\to B\) is surjective. Then, there is a function 
  \(g: B \to A\) so that \(f \circ g = \one_{B}\).

  Suppose \(Z\) is any set and \(\alpha',\alpha'' \colon B \to Z\) are any pair of 
  functions so that \(\alpha' \circ f = \alpha'' \circ f\).

  This implies that 

  \[
    \alpha' \circ (f \circ g) = \alpha'' \circ (f \circ g)
  \]

  from which we see that 

  \[
    \alpha' \circ \one_{B} = \alpha'' \circ \one_{B}
  \]

  from which we derive that \(\alpha' = \alpha''\). This completes one side of the proof. 

  \((\Longleftarrow)\)
  Now, suppose \(f\) is an epimorphism. Set \(Z \coloneqq \{0,1\}\), \(\alpha' \coloneqq \one_{B}\) 
  and \(\alpha'' \coloneqq \one_{\image f}\).
  Notice that \(\alpha'\) and \(\alpha''\) agree on the image of \(f\). Therefore, by the 
  fact that \(f\) is an epimorphism, they are identical. Since the indicators of two sets 
  are identical, the two sets are equal and this concludes our proof.
\end{proof}

\begin{definition}
  A category \(\catC\) consists of 
  \begin{itemize}

    \item A class \(\Obj(\catC)\) of \emph{objects} of the category and

    \item for every pair of objects \(A,B\) of \(\catC\), a set \(\HomC(A,B)\)
      of \emph{morphisms}, with the following properties: 

      \begin{enumerate}

        \item For every object \(A\) of \(\catC\) (at least) one morphism
          \(1_{A}\in \HomC(A, A)\), called the \emph{identity} on \(A\).

        \item Any two morphisms \(f\in \HomC(A,B)\) and 
          \(g\in \HomC(B,C)\) determine another morphism, denoted 
          \(gf\in \HomC(A,C)\). That is, for every triple of objects \(A,B\) and 
          \(C\) of \(\catC\), there is a function (of sets)

          \[
            \HomC(A,B) \times \HomC(B,C) \longrightarrow \HomC(A, C)
          \]

          and the image of the pair \(f,g\) is denoted \(gf\).

        \item The previous composition law is associative: if \(f\in \HomC(A,B)\), 
          \(g\in \HomC(B,C)\) and \(h\in \HomC(C, D)\), then 

          \[
            (hg)f = h(gf)
          \]

        \item The identity morphisms are identities with respect to composition, that is,
          for all \(f\in \HomC(A,B)\) we have 

          \[
            f1_{A} = f, \qquad 1_{B}f = f
          \]

      \end{enumerate}

      A further requirement is that \(\HomC(A,B)\) and \(\HomC(C,D)\) 
      have to be disjoint unless \(A = C\) and \(B = D\).
  \end{itemize}
  
\end{definition}

\begin{notation}[Endomorphism]
  A morphism \(f\) of an object \(A\in \Obj(\catC)\) of a category \(\catC\) to itself is 
  called an \emph{endomorphism}; \(\HomC(A,A)\) is denoted by \(\End_{\catC}(A)\)
\end{notation}



\printbibliography
\end{document}
